\scanpage{23}\mspage{8}
où $\delta$ est le discriminant du second membre de (3.1). Le discriminant $\delta$ étant invariant par translation, il suffit de vérifier cette identité algébrique pour $3$ inversible et $\mathrm{a}_{2}=0$.

Lorsque $\mathrm{c}_{4}$ est inversible, un et un seul choix de $\mathrm{x}$ et $\mathrm{y}$ donne $\mathrm{d}_{6}=0$ ? i.e. $\mathrm{a}_{2}\,\mathrm{a}_{4}-9\,\mathrm{a}_{6}=0$. On a, lorsque cette condition est remplie,
\[
-\mathrm{c}_{6}=2^{6}(\mathrm{a}_{2}^{3}-3^{3}\,\mathrm{a}_{6})\qquad
\Delta=2^{6}(-\mathrm{a}_{2}^{3}\,\mathrm{a}_{6}-\mathrm{a}_{4}^{3}+2.3^{3}\,\mathrm{a}_{6}^{2})
\tag*{(3.3)}
\]
et
\[
\begin{aligned}[t]
\mathrm{a}_{2}\,\mathrm{c}_{4}&=-2^{-2}\,\mathrm{c}_{6}\\
\mathrm{a}_{4}\,\mathrm{c}_{4}^{2}&=-2^{2}\,3^{2}\,\Delta\\
\mathrm{a}_{6}\,\mathrm{c}_{4}^{3}&=\mathrm{c}_{6}\,\Delta\qquad.
\end{aligned}
\tag*{(3.4)}
\]

Il suffit de vérifier ces identités en caractéristique $0$, auquel cas la troisième est le produit des deux premières, puisque $\mathrm{d}_{6}=0$.

\uline{Proposition} 3.5.\quad \uline{Au dessus de} $\operatorname{Spec}(\handwritten{\mathbb{Z}}[2^{-1}])$, \uline{il existe un schéma de modul}\tikz[baseline=(ending.base)]{\node[inner sep=0pt] (ending) {\uline{e\handwritten{s}}};\handwritten{\draw[overlay,rotate around={-6:(ending.center)}] (ending.center) ellipse [x radius=.6em,y radius=1.35em];}}\uline{ pour les courbes de genre un munies d'une forme différentielle invariante inversible et telles que} $\mathrm{c}_{4}$ \uline{soit inversible. Ce schéma} \reversemarginpar\margin{\handwritten{\tikz[baseline=-.15ex]{\draw[thin,gray] (0,0) -- (.63em,.14em);}}}\normalmarginpar\uline{est}
\[
\operatorname{Spec}(\handwritten{\mathbb{Z}}[2^{-1}][\mathrm{a}_{2},\,\mathrm{a}_{4},\,\mathrm{a}_{6},\,(\mathrm{a}_{2}^{2}-3\,\mathrm{a}_{4})^{-1}]/(\mathrm{a}_{2}\,\mathrm{a}_{4}-9\,\mathrm{a}_{6}))
\]
\uline{avec pour courbe universelle} (\uline{en coordonnées non homogè}\tikz[baseline=(ending.base)]{\node[inner sep=0pt] (ending) {\uline{ne}\handwritten{s})};\handwritten{\draw[overlay] ([xshift=1.3em,yshift=2.16em]ending.center) .. controls +(-.8em,.42em) and +(.3em,1.1em) .. ([xshift=-.63em,yshift=.2em]ending.center) .. controls +(-.72em,-1.6em) and +(-.76em,-1.6em) .. ([xshift=1.35em,yshift=-.65em]ending.center) .. controls +(.7em,.8em) and +(.5em,1.6em) .. ([xshift=1.32em,yshift=2.22em]ending.center);}}
\[
\mathrm{y}^{2}=\mathrm{x}^{3}+\mathrm{a}_{2}\,\mathrm{x}^{2}+\mathrm{a}_{4}\,\mathrm{x}+\mathrm{a}_{6}
\]
\clearpage
% archive sheet number: 22
